\documentclass[../../../include/open-logic-section]{subfiles}

\begin{document}

\olfileid{sth}{spine}{zf}
\olsection{$\Z$ 与 $\ZF$：一个里程碑}

引入基础公理后，我们达到了另一个重要的里程碑。我们已经考察了理论 $\Zminus$ 和 $\ZFminus$，我们曾说它们是“减去”了某个东西的理论。而这个东西正是基础公理。因此：

\begin{defn}
理论 $\Z$ 是在 $\Zminus$ 中加入基础公理。因此其公理为：外延公理、并集公理、对集公理、幂集公理、无穷公理、基础公理，以及分离公理模式的所有实例。

理论 $\ZF$ 是在 $\ZFminus$ 中加入基础公理。换言之，$\ZF$ 是在~$\Z$ 中加入替换公理模式的所有实例。
\end{defn}

不过，你可能会有一个疑问：如果正则公理在 $\ZFminus$ 中与基础公理等价，且正则公理的合理性是显然的，那为什么要舍近求远，将基础公理作为我们的基本公理，而不是正则公理呢？

撇开历史原因（涉及谁在何时提出了何种表述）不谈，根本原因在于，基础公理的表述无需依赖~$V_\alpha$ 的定义。这一定义依赖于 \olref[recursion]{sec} 的全部工作：我们需要证明超限递归定理，以表明该定义是合理的。但我们证明超限递归定理时用到了\emph{替换公理}。因此，尽管基础公理与正则公理相对于 $\ZFminus$ 是等价的，它们相对于 $\Zminus$ 却不等价。

事实上，情况比这简单的一提所暗示的更为严峻。尽管这远远超出了本书的范围，但事实证明，$\Zminus$ 和 $\Z$ 都太弱，以至于无法定义 $V_\alpha$ 层级。所以，如果你只在 $\Z$ 中工作，那么正则公理（按我们此前的表述方式）甚至连\emph{意义}都没有。这就是为什么我们正式采纳的公理是基础公理，而不是正则公理。

从现在起，我们将在 $\ZF$ 中工作（除非另有说明），不再特别声明。

\end{document}